% ============================================================
% Chapter 3 — System Design (target 11–13 pp)
% Sources: carla_multi_agent_env.py, route_planner.py,
%   curriculum_batch_manager.py, centralized_critic.py,
%   levels.yaml, curriculum_batch.yaml, train_mappo.yaml.
% Status: FULL DRAFT 2026-07-14 (§3.1–3.7 + Tab. 1–2 + Fig. 1 TikZ).
% Pending: Fig. 2 (Town03 screenshot), Fig. 3 (route examples),
%          Fig. 4 (curriculum state machine) — see \todo markers.
% ============================================================
\chapter{System Design}
\label{ch:system}

This chapter describes the experimental platform: a multi-agent
reinforcement-learning environment built on the CARLA simulator, in which
three vehicles and three pedestrians are trained jointly with MAPPO under
two alternative training regimes --- curriculum and mixed-batch --- that
share every other component of the stack. The design goal is that the
training regime is the \emph{only} independent variable of the main
experiment; everything documented in this chapter is common to both arms.

% ------------------------------------------------------------
\section{Architecture Overview}
\label{sec:architecture}

The platform consists of four layers (\cref{fig:architecture}): the CARLA
server, a custom multi-agent environment, the RLlib training stack, and a
logging/evaluation pipeline.

\paragraph{Simulator.}
CARLA 0.9.16~\cite{dosovitskiy2017carla}, built on Unreal Engine~4, runs as
a standalone server. The environment drives it in \emph{synchronous mode}
with a fixed physics step of \SI{0.05}{s} (\SI{20}{Hz}): the simulator
advances only when the environment issues a tick, which makes experience
collection deterministic with respect to scheduling and independent of
wall-clock speed. Scripted (non-learning) traffic --- 15 NPC vehicles and
30 NPC pedestrians in every training level --- is managed by the CARLA
Traffic Manager, also in synchronous mode. Training uses the urban map
\emph{Town03}; \emph{Town05} is reserved exclusively for the
generalization test and is never seen during training.

\paragraph{Multi-agent environment.}
A custom environment (\texttt{CarlaMultiAgentEnv}) exposes the simulation
to RLlib as a multi-agent interface with six learning agents: three
vehicles and three pedestrians (walkers). At every reset the environment
spawns the agents, assigns each an individual route (\cref{sec:routes}),
and configures the scenario according to the difficulty level selected by
the curriculum/batch manager (\cref{sec:regimes}). An episode lasts at most
1{,}000 control steps (\SI{50}{s} of simulated time at \SI{20}{Hz}).

\paragraph{Learning stack.}
Training uses MAPPO~\cite{yu2022surprising} as implemented on
Ray/RLlib~2.10.0 with PyTorch~2.7 (Python 3.11.9), in the
\emph{centralized training with decentralized execution} (CTDE) paradigm.
Two policies are learned: a \emph{vehicle policy} shared by the three
vehicles and a \emph{pedestrian policy} shared by the three pedestrians
(parameter sharing within each class, no sharing across classes, since the
two classes have different observation and action spaces). Each actor is a
multi-layer perceptron with two hidden layers of 256 units. A centralized
critic receives a 225-dimensional global observation that concatenates
per-agent information (\cref{sec:obs-actions}); the critic encoder is an
MLP by default, with a GNN variant used as a secondary experimental axis
(\cref{sec:critic-variants}).

\paragraph{Logging and evaluation.}
At the end of every episode the training process appends one JSON record
\emph{per agent} (six per episode) to an \texttt{episodes.jsonl} file:
termination reason, route completion, path efficiency, speed, and
route-generation diagnostics. All results in \cref{ch:results} are computed
from these records --- either the training stream (cumulative, on-policy
metrics) or the records produced by a separate evaluation pipeline, which
reloads a trained checkpoint and runs four fixed scenarios (easy, medium,
hard on Town03; test on Town05) of 100 episodes each
(\cref{sec:critic-variants}).

\begin{figure}[t]
  \centering
  \begin{tikzpicture}[
      node distance=7mm and 12mm,
      box/.style={draw, rounded corners=2pt, align=center,
                  minimum width=52mm, minimum height=11mm, font=\small},
      lbl/.style={font=\scriptsize, midway},
      >={Stealth[length=2.5mm]}]
    \node[box] (carla) {CARLA 0.9.16 server\\ \scriptsize Town03 (train) / Town05 (test), TM};
    \node[box, below=of carla] (env) {\texttt{CarlaMultiAgentEnv}\\ \scriptsize 3 vehicles + 3 walkers, route planner};
    \node[box, right=of env] (mgr) {Curriculum / batch\\ manager \scriptsize (level per episode)};
    \node[box, below=of env] (rllib) {Ray/RLlib 2.10 --- MAPPO (CTDE)\\ \scriptsize vehicle policy, pedestrian policy, centralized critic};
    \node[box, below=of rllib] (log) {\texttt{episodes.jsonl}\\ \scriptsize 6 records/episode};
    \node[box, right=of log] (eval) {Evaluation pipeline\\ \scriptsize $4$ scenarios $\times$ 100 ep};
    \draw[->] (env) -- node[lbl, right] {tick / control} (carla);
    \draw[<-] ([xshift=-8mm]env.north) -- node[lbl, left] {sensors, state} ([xshift=-8mm]carla.south);
    \draw[->] (mgr) -- node[lbl, above] {level} (env);
    \draw[->] (env) -- node[lbl, right] {obs, rewards} (rllib);
    \draw[<-] ([xshift=-8mm]env.south) -- node[lbl, left] {actions} ([xshift=-8mm]rllib.north);
    \draw[->] (rllib) -- (log);
    \draw[->] (rllib) -| node[lbl, above, pos=0.25] {checkpoint} (eval);
  \end{tikzpicture}
  \caption{Component architecture. The curriculum/batch manager is the only
  component that differs between the two experimental arms, and only in its
  level-selection rule (\cref{sec:regimes}).}
  \label{fig:architecture}
\end{figure}

% ------------------------------------------------------------
\section{A Shared-World Multi-Agent Environment}
\label{sec:shared-world}

\subsection{Parallel task coverage}
\label{sec:parallel-coverage}

The six learning agents do not cooperate toward a joint goal: each agent is
assigned its \emph{own} route and is evaluated individually on completing
it. Because all six routes are placed in the same living world, a single
episode simultaneously exercises six different micro-tasks --- different
road geometries, different interactions with scripted traffic, and, most
importantly, mutual interactions between learning agents (e.g.\ a learning
vehicle yielding to a learning pedestrian). We refer to this design as
\emph{parallel task coverage}: per wall-clock episode, the system collects
experience across six task instances rather than one.

This framing comes with two caveats that the analysis must respect. First,
the six per-episode samples are \emph{correlated} --- the agents share the
world state, so a traffic jam or a blocked crosswalk affects several
records at once. Second, from the perspective of any single agent the
environment is \emph{non-stationary} during training, because the other
five agents' policies change over time; this is the standard motivation for
the CTDE critic (\cref{ch:background}). Parallel task coverage is a
property of the platform, present identically in both training regimes: it
is part of the experimental frame, not an experimental variable.

\subsection{Per-agent routes and seeding}
\label{sec:seeding}

At every reset, each agent receives an independent route whose target
length is set by the active difficulty level (30/60/\SI{100}{m};
\cref{sec:routes}). Route generation is deterministically seeded per agent
with a seed sequence built from
$\langle\text{traffic seed},\ \text{reset counter},\ \mathrm{crc32}(\text{agent
id})\rangle$, so that (i) two training runs launched with the same
experiment seed on the same code see the same sequence of route-generation
draws, and (ii) route draws are decoupled across agents. This makes paired
A/B comparisons route-matched by construction, a property the experimental
protocol relies on (\cref{ch:methodology}).

\subsection{Termination taxonomy and per-agent records}
\label{sec:termination}

Each agent terminates its episode independently with exactly one of six
reasons, classified by a simulator-independent rule
(\texttt{episode\_classification.py}):

\begin{itemize}
  \item \texttt{route\_complete} --- the agent reached the end of its
        route. This is the \emph{only} outcome counted as success.
  \item \texttt{route\_short} --- the agent reached the end of a route
        whose optimal length turned out to be shorter than the level target
        by more than a tolerance factor (degenerate route-generation
        fallback). Reaching it is \emph{not} counted as success: this guard
        prevents success-rate inflation through trivially short routes.
  \item \texttt{collision} --- contact with any actor or static geometry.
  \item \texttt{offroad} --- the vehicle moved more than \SI{5}{m} away
        from the nearest driving-lane centre (vehicles only).
  \item \texttt{stuck} --- prolonged lack of progress along the route.
  \item \texttt{timeout} --- the 1{,}000-step episode cap elapsed first.
\end{itemize}

Every episode therefore yields exactly six agent-level records (three
vehicle, three pedestrian) in \texttt{episodes.jsonl}. All aggregate
metrics in this thesis are computed at the agent level over these records,
always reported separately for vehicles and pedestrians; the measurement
rules, including deduplication and the strict success criterion, are
specified in \cref{ch:methodology}.

% ------------------------------------------------------------
% Sources (verified 2026-07-14): carla_multi_agent_env.py
%   _get_vehicle_obs L1291-1381, _get_pedestrian_obs L1383-1445,
%   _fill_nearby L1740-1764, constants L84-88; global obs layout
%   centralized_critic.py (compute_global_obs_dim_with_mask, L115).
\section{Observation and Action Spaces}
\label{sec:obs-actions}

Each agent receives a compact, class-specific observation vector with all
components normalized to $[-1, 1]$ and sanitized against non-finite values
before being returned to the learner. Vehicles observe 47 dimensions and
pedestrians 26; \cref{tab:obs} lists both layouts. Three design choices
deserve comment.

\paragraph{Route-aware features.}
Beyond the classic ego-kinematics and nearest-neighbour features, both
classes observe a short \emph{preview} of their own route (relative
offsets of upcoming waypoints in the ego frame) together with heading
error and, for vehicles, upcoming turn angle and signed lateral error.
These features give the policy the local geometry of the assigned task
without revealing anything about the other agents' goals.

\paragraph{Hazard preview channel.}
A shared routine estimates, along the agent's forward corridor, a
time-to-collision risk and a corridor-occupancy fraction, computed
separately against vehicles (corridor \SI{3.5}{m} wide, \SI{40}{m} ahead,
\SI{5}{s} horizon) and against pedestrians (\SI{3.0}{m}, \SI{30}{m},
\SI{4}{s}); pedestrians observe the vehicle channel only (\SI{4.0}{m},
\SI{30}{m}, \SI{4}{s}). The same four vehicle-side risk values are reused
by the reward function as a safety gate (\cref{sec:rewards}), so the
quantity that modulates the reward is always visible to the policy.

\paragraph{Markov-consistency features.}
The last three vehicle dimensions expose internal environment state that
the reward depends on: the number of steps without waypoint progress
(normalized by 300), the state of the anti-loop penalty, and the
normalized time remaining before episode truncation. Without these, the
penalties of \cref{sec:rewards} and the fixed horizon would act on
information the agent cannot observe --- a state-aliasing defect that was
identified and corrected during development (\cref{ch:methodology}).

\begin{table}[t]
  \centering\small
  \caption{Observation layouts. Slashes denote normalization constants;
  neighbour features are expressed in the ego frame.}
  \label{tab:obs}
  \begin{tabular}{@{}llp{8.0cm}@{}}
    \toprule
    Dims & Group & Content \\
    \midrule
    \multicolumn{3}{@{}l}{\emph{Vehicle (47D)}} \\
    0--8   & Ego kinematics   & velocity (/30), acceleration (/10), speed (/120 km/h), planar heading \\
    9--12  & Route tracking   & distance (/\SI{50}{m}) and bearing (/$\pi$) to next waypoint; signed in-lane lateral offset (lane widths); continuous route progress \\
    13--14 & Traffic light    & at-light flag; state (red 1.0, yellow 0.5, green 0.0) \\
    15--23 & Nearby vehicles  & 3 nearest within \SI{50}{m}: longitudinal and lateral offset (/50), relative speed (/30) \\
    24     & Control memory   & previous steering command \\
    25--30 & Route preview    & waypoints $+1/+3/+5$: longitudinal and lateral offsets (/50) \\
    31--33 & Route geometry   & heading error; upcoming turn angle; signed lateral error to route \\
    34--39 & Nearby pedestrians & 2 nearest within \SI{50}{m}, same 3 features \\
    40--43 & Hazard preview   & time-to-collision risk and corridor occupancy, vs.\ vehicles and vs.\ pedestrians \\
    44--46 & Progress/time    & steps without progress (/300); loop-penalty flag; time remaining \\
    \midrule
    \multicolumn{3}{@{}l}{\emph{Pedestrian (26D)}} \\
    0--2   & Position         & world position ($x, y$ /500; $z$ /50) \\
    3--6   & Kinematics       & velocity (/5), speed (/5 m/s) \\
    7--8   & Heading          & planar forward vector \\
    9--10  & Waypoint         & distance (/\SI{100}{m}) and bearing to current waypoint \\
    11     & Surface          & on-sidewalk flag \\
    12--17 & Nearby vehicles  & 2 nearest within \SI{50}{m}: same 3 features \\
    18     & Progress         & route completion fraction \\
    19--22 & Route preview    & waypoints $+1/+3$: longitudinal and lateral offsets (/100) \\
    23     & Route geometry   & heading error \\
    24--25 & Hazard preview   & vehicle time-to-collision risk and corridor occupancy \\
    \bottomrule
  \end{tabular}
\end{table}

\paragraph{Action spaces.}
Both classes use continuous two-dimensional actions. The vehicle action is
$[a_{\mathrm{tb}}, a_{\mathrm{steer}}] \in [-1,1]^2$, where positive
$a_{\mathrm{tb}}$ maps to throttle and negative to brake (reverse gear is
disabled), and $a_{\mathrm{steer}}$ is the steering command. The
pedestrian action is $[a_{\mathrm{speed}}, a_{\Delta\mathrm{h}}]$ with
$a_{\mathrm{speed}} \in [0,1]$ scaling a maximum walking speed of
\SI{5.0}{m/s} and $a_{\Delta\mathrm{h}} \in [-1,1]$ commanding a heading
change of up to $\pm 45^{\circ}$ per step.

\paragraph{Global observation for the centralized critic.}
During training only, the critic receives a fixed-slot concatenation of
all agents' observations followed by a per-agent alive mask:
$[\,v_0\,|\,v_1\,|\,v_2\,|\,p_0\,|\,p_1\,|\,p_2\,|\,\text{alive}_{6}\,]$,
i.e.\ $3 \times 47 + 3 \times 26 + 6 = 225$ dimensions. The alive mask is
necessary because agents terminate independently
(\cref{sec:termination}): slots of terminated agents are masked rather
than silently reused. Actors never see this vector, preserving
decentralized execution.

% ------------------------------------------------------------
% Sources (verified 2026-07-14): carla_multi_agent_env.py
%   _vehicle_reward L1866-1960 ("Vehicle Reward v5"),
%   _pedestrian_reward L1962-2010 ("Pedestrian Reward v5").
\section{Reward Design}
\label{sec:rewards}

Rewards are per-agent and identical in both training regimes.
\Cref{tab:rewards} lists every component with its trigger and magnitude;
this section explains the structure. The final form is the outcome of the
gate-driven iterations documented in \cref{sec:iterations}; no claim of
optimality is attached to the individual coefficients.

\paragraph{Progress core.}
For both classes the dominant signal is a sparse bonus per route waypoint
reached ($+100$ for vehicles, $+50$ for pedestrians), complemented by a
dense shaping term proportional to the per-step reduction of the distance
to the next waypoint. Waypoints skipped by cutting across the route are
penalized, so progress cannot be gamed by shortcuts.

\paragraph{Safety terms.}
A collision costs vehicles $-500$, applied once at the step of first
contact (the episode then terminates). Two further vehicle terms penalize
leaving the lane (linear beyond \SI{2}{m} from the lane centre; episode
termination occurs at \SI{5}{m}, \cref{sec:termination}) and exceeding
50~km/h. Pedestrian collisions cost $-50$, and pedestrians receive a
small per-step incentive to stay on sidewalks rather than on driving
lanes.

\paragraph{Hazard-gated anti-stall shaping (vehicles).}
The empirically most delicate block counteracts the \emph{defensive
equilibrium} in which a vehicle stops indefinitely to avoid any risk. All
of its terms are suspended while the vehicle is legitimately waiting at a
red or yellow light. The scalar $\mathit{hazard} =
\max(\text{TTC}_{\mathrm{veh}}, \text{occ}_{\mathrm{veh}},
\text{TTC}_{\mathrm{ped}}, \text{occ}_{\mathrm{ped}})$ summarizes the
hazard-preview channel of \cref{sec:obs-actions}, and a boolean
$\mathit{safe\_to\_push} = (\mathit{hazard} < 0.85)$ gates the pushing
terms: an idle penalty below 2.5~km/h; a start bonus proportional to
throttle, route alignment, and an urgency factor that grows with the time
spent without progress; and a minimum-speed shaping around a target of
8~km/h, active only when the vehicle is aligned with its route (alignment
$> 0.25$). Above the hazard threshold every pushing incentive disappears,
so the policy is never rewarded for accelerating into a risky corridor.
Independently of the gate, a no-progress ramp (after 100 steps without
reaching a waypoint) and an anti-loop penalty make prolonged stalling and
circling strictly unprofitable, and a steering-smoothness pair discourages
jerky control.

\paragraph{Pedestrian pace band.}
Pedestrians earn a per-step bonus while walking within a comfort band of
$[1.2, 2.6]$~m/s, with an idle penalty below 0.15~m/s and a running
penalty above 3.0~m/s. The band is intentionally wide: at 2.0~m/s a
pedestrian covers \SI{100}{m} --- the hard-level target --- within the
\SI{50}{s} horizon. The bonus-only (asymmetric) design of this band and
its interaction with the vehicles' hazard gate are analysed empirically in
\cref{ch:results}.

\begin{table}[t]
  \centering\small
  \caption{Reward components. ``Gated'' vehicle rows are suspended while
  waiting at a red/yellow light; rows marked $^{\dagger}$ additionally
  require $\mathit{safe\_to\_push}$ ($\mathit{hazard} < 0.85$) and route
  alignment $> 0.25$.}
  \label{tab:rewards}
  \begin{tabular}{@{}p{3.4cm}p{5.6cm}p{4.6cm}@{}}
    \toprule
    Component & Trigger & Value \\
    \midrule
    \multicolumn{3}{@{}l}{\emph{Vehicle}} \\
    Waypoint reached      & per waypoint advanced & $+100$ \\
    Skipped waypoint      & per waypoint skipped & $-10$ \\
    Dense progress        & change in distance to next waypoint & $+4.0$ per metre closed \\
    Collision             & first contact (terminal) & $-500$ \\
    Off-lane              & $>$ \SI{2}{m} from lane centre & $-0.5$ per metre \\
    Overspeed             & speed $> 50$ km/h & $-0.10$ per km/h over \\
    Idle (gated)          & speed $< 2.5$ km/h & $-0.15$ \\
    Start bonus (gated$^{\dagger}$) & throttle while idle & $(0.20 + 0.30\,u)\cdot\text{thr}\cdot\text{align}$, brake symmetric; $u = \min(\text{no-progress}/200, 1)$ \\
    Min-speed shaping (gated$^{\dagger}$) & speed below / above 8 km/h & $-0.04\,(8 - v)\cdot\text{align}$ / $+0.15\cdot\text{align}$ \\
    No-progress ramp (gated) & $>$ 100 steps without waypoint & $-0.004$ per step, capped $-1.0$ \\
    Smooth steering       & $|\Delta \text{steer}| < 0.1$ at $> 5$ km/h & $+0.1$ \\
    Steering jerk         & $|\Delta \text{steer}| > 0.5$ & $-0.3$ \\
    Anti-loop             & loop detector active & $-1.0$ \\
    \midrule
    \multicolumn{3}{@{}l}{\emph{Pedestrian}} \\
    Waypoint reached      & per waypoint advanced & $+50$ \\
    Skipped waypoint      & per waypoint skipped & $-10$ \\
    Dense progress        & change in distance to waypoint & $+2.0$ per metre closed \\
    Collision             & first contact (terminal) & $-50$ \\
    Sidewalk / roadway    & per step on sidewalk / driving lane & $+0.2$ / $-0.3$ \\
    Walking pace          & speed $\in [1.2, 2.6]$ m/s & $+0.3$ \\
    Idle                  & speed $< 0.15$ m/s & $-0.15$ \\
    Running               & speed $> 3.0$ m/s & $-0.7$ per m/s over \\
    Anti-loop             & loop detector active & $-1.0$ \\
    \bottomrule
  \end{tabular}
\end{table}

% ------------------------------------------------------------
% Sources (verified 2026-07-14): route_planner.py (min_route_ratio=0.5,
%   max_route_ratio=2.0 L145-146/206-207, max_cross=3 L261/329, short-chain
%   rejection L368-369); levels.yaml (levels_path 30/60/100, test Town05 80);
%   route_source labels carla_multi_agent_env.py L1040-1159.
\section{Route Generation and Difficulty Levels}
\label{sec:routes}

Difficulty is defined by the \emph{target length} of the per-agent routes,
with all other scenario parameters held constant: the three training
levels (easy, medium, hard) place both vehicle and pedestrian targets at
30, 60 and \SI{100}{m} respectively, always on Town03 and always with the
same scripted traffic (15 NPC vehicles, 30 NPC pedestrians). A fourth,
evaluation-only level (\emph{test}) uses Town05 with \SI{80}{m} targets
and the same traffic counts. The configuration also defines level families
where difficulty varies by traffic density instead of route length; these
are not used in this thesis (\cref{sec:campaign}).

\paragraph{Vehicle routes.}
Vehicle routes are planned on the road topology with CARLA's global route
planner: candidate destinations are sampled around the target distance and
a candidate is accepted only if its optimal path length lies within
$[0.5, 2.0] \times$ the target, with up to 32 attempts per reset. If no
candidate survives, the environment falls back to a legacy waypoint-chain
generator. The applied generator is recorded per record
(\texttt{route\_source} $\in$ \{\texttt{grp},
\texttt{legacy\_fallback}\}), so the share of fallback routes is always
measurable.

\paragraph{Pedestrian routes.}
Pedestrian routes are chains of sidewalk segments connected through
crosswalks. Two constraints shape them: chains whose cumulative length
falls below $0.5 \times$ the target are rejected (the same lower-bound
guard that backs the \texttt{route\_short} demotion of
\cref{sec:termination}), and the number of crosswalk crossings per route
is capped at three --- a bound on the pedestrians' on-roadway exposure
whose value was set empirically (\cref{sec:iterations}). The generator
records its outcome per record (\texttt{sidewalk\_crosswalk},
\texttt{sidewalk\_distance}, \texttt{respawn}, \texttt{legacy\_chain}),
together with feasibility diagnostics (optimal length, target error,
under-target and too-short flags, candidate attempts). Pedestrian route
feasibility on Town03 is a structural limitation of the platform ---
sidewalk topology often cannot support the nominal medium/hard targets ---
and is analysed separately in \cref{sec:ped-limits}, where environment
limits are explicitly distinguished from policy limits.

\todo{Fig.~3: example routes per level. Generator script ready at
\texttt{docs/thesis/scripts/make\_fig3\_route\_examples.py} (requires a
live CARLA server; user runs it); then include
\texttt{figures/route\_examples.pdf} here.}

% ------------------------------------------------------------
% Sources (verified 2026-07-14): curriculum_batch.yaml (all thresholds);
%   curriculum_batch_manager.py (_balanced_policy_success_details L391-423:
%   min over policies; window_full L129/382; BatchLevelSampler L902-940).
\section{Curriculum and Mixed-Batch Training Regimes}
\label{sec:regimes}

The training regime is the independent variable of the main experiment.
Both regimes are implemented by the same manager component, consulted once
per episode reset to select the difficulty level; they share the level
definitions of \cref{sec:routes}, the total budget of $3 \times 10^6$
environment steps, and every other component described in this chapter.
They differ \emph{only} in the level-selection rule.

\subsection{Curriculum arm: competence-gated progression}
\label{sec:curriculum-arm}

Training starts with only the easy level unlocked. Higher levels unlock
through a competence gate evaluated on a sliding window of the last 50
episodes: the windowed success rate is computed separately for the vehicle
policy and the pedestrian policy, and the gate binds on the \emph{minimum}
of the two --- progression requires the weaker class to be competent, not
the average. Medium unlocks at a windowed success rate of at least 70\%,
hard at 60\%; the window must be full, and a minimum share of the total
budget must already have been spent at the preceding levels (20\% on easy
before medium; 28\% on medium before hard). Two design decisions deserve
justification:

\begin{itemize}
  \item \textbf{Windowed, not cumulative.} The cumulative success rate is
        a lagging integrator dragged down by cold-start failures: a policy
        can be competent \emph{now} while its since-birth average is still
        far below threshold. Gating on a sliding window measures current
        competence.
  \item \textbf{Pressure caps.} If a threshold is never met, the level
        force-unlocks anyway once a fixed share of the total budget has
        been consumed (30\% for medium, 70\% for hard). This guarantees
        that every level is visited within the budget, so the curriculum
        arm can never terminally stall on easy --- at worst it degrades
        toward a delayed mixed exposure.
\end{itemize}

Once unlocked, levels are sampled with weights that normalize to
$0.30/0.35/0.35$ (easy/medium/hard) under budget constraints (easy at most
30\% of the budget once medium is available, medium at most 35\%, hard at
least 35\%). Newly unlocked levels pass through a short \emph{probation}
phase with reduced sampling pressure (hard is temporarily de-emphasized at
weight 0.85), smoothing the transition instead of switching abruptly.

\subsection{Batch arm: stratified interleaving}
\label{sec:batch-arm}

The mixed-batch arm is the \emph{no-ordering control}: all three levels
are available from the first episode. Levels are sampled by stratified
shuffling without replacement with a window of three: each consecutive
window of three episodes drawn from one sampler instance is a uniformly
random permutation of \{easy, medium, hard\}, giving uniform $1/3$
exposure with a maximum imbalance of one episode per sampler --- avoiding
both the drift of i.i.d.\ sampling and any fixed deterministic ordering
within the window. In the realized campaign logs, where episodes from
parallel rollout workers interleave, the global level shares stay within
$\approx 2$ percentage points of uniform and all three levels appear
within the first twenty episodes (\cref{sec:integrity}).

\subsection{What the comparison isolates}
\label{sec:isolation}

Under this construction, the two arms consume the same budget over the
same task distribution in expectation; what differs is \emph{when} each
part of that distribution is presented and whether presentation is
conditioned on measured competence. Any behavioural difference between the
arms is therefore attributable to the ordering policy, up to run-to-run
variance --- which is quantified separately in \cref{ch:results} and
discussed in \cref{sec:threats}.

\begin{figure}[t]
  \centering
  \begin{tikzpicture}[
      state/.style={draw, rounded corners=2pt, align=center,
                    minimum width=34mm, minimum height=13mm, font=\small},
      gate/.style={font=\scriptsize, align=center},
      >={Stealth[length=2.5mm]}, node distance=17mm]
    \node[state] (s1) {easy only\\ \scriptsize sampling: easy 1.0};
    \node[state, right=of s1] (s2) {easy $+$ medium\\ \scriptsize probation, then\\ \scriptsize $\propto 0.30/0.35$};
    \node[state, right=of s2] (s3) {easy $+$ medium $+$ hard\\ \scriptsize probation (hard $\times$0.85), then\\ \scriptsize $\propto 0.30/0.35/0.35$};
    \draw[->] (s1.north east) to[bend left=28]
      node[gate, above] {win.\ SR $\geq 0.70$\\ $\wedge$ easy share $\geq 0.20$} (s2.north west);
    \draw[->, dashed] (s1.south east) to[bend right=28]
      node[gate, below] {cap: 30\% budget} (s2.south west);
    \draw[->] (s2.north east) to[bend left=28]
      node[gate, above] {win.\ SR $\geq 0.60$\\ $\wedge$ med.\ share $\geq 0.28$} (s3.north west);
    \draw[->, dashed] (s2.south east) to[bend right=28]
      node[gate, below] {cap: 70\% budget} (s3.south west);
  \end{tikzpicture}
  \caption{Curriculum progression. Solid edges: competence gates on the
  windowed success rate (minimum over the two policies, window of 50
  episodes, window full) with budget-share floors. Dashed edges: pressure
  caps guaranteeing progression within the budget. In the campaign's
  MLP-curriculum run (seed 999), medium unlocked at episode 842 of 3{,}220
  (26.1\%) and hard at episode 2{,}286 (71.0\%), for a final level mix of
  32.5\% easy / 45.8\% medium / 21.6\% hard --- timings consistent with
  the pre-campaign trunk replicates.}
  \label{fig:curriculum-fsm}
\end{figure}

% ------------------------------------------------------------
% Sources (verified 2026-07-14): centralized_critic.py (GraphConvLayer
%   L336-356, GNNCriticEncoder L414-497, critic MLP 225->256->256->1);
%   run_config.json of campaign runs (all optimizer values, byte-identical
%   across arms except mode/use_gnn).
\section{Critic Encoder Variants and Training Pipeline}
\label{sec:critic-variants}

\subsection{Critic encoders}
\label{sec:critic-encoders}

In the default configuration the centralized critic is an MLP mapping the
225-dimensional global observation through two hidden layers of 256 units
to a scalar value. The secondary experimental axis
(\cref{sec:campaign}) replaces this encoder with a graph-based one, while
\emph{both actors remain the identical MLPs} of \cref{sec:architecture}:
any behavioural difference between the two variants is therefore
attributable to value estimation during training, not to policy capacity,
and at execution time the two variants are indistinguishable in
architecture.

The GNN encoder splits the global observation back into the six per-agent
slots, projects each through a per-class linear layer into a
64-dimensional node embedding, and applies two rounds of message passing
with GraphSAGE-style mean aggregation~\cite{hamilton2017graphsage}: each
node combines its own embedding with the mean of the other \emph{alive}
agents' embeddings through separate learned weights, followed by an ELU
activation. A masked mean-pool over alive nodes produces the input to the
value head. The alive mask from the global observation drives both the
aggregation and the pooling, so terminated agents drop out of the
critic's summary instead of contributing stale slots. (The codebase also
provides an attention-based critic encoder and PopArt value
normalization behind configuration flags; both remain disabled throughout
this thesis.)

\subsection{Optimization settings}
\label{sec:optimization}

All campaign runs share the same optimizer configuration, byte-identical
across arms except for the regime and encoder flags: learning rate
$5\times10^{-4}$; discount $\gamma = 0.997$; GAE $\lambda = 0.95$; PPO
clip 0.2 with an additional KL penalty (target 0.02); gradient clipping at
0.5; value-loss coefficient 0.05 with the value clip effectively disabled
($10^6$); training batch 8{,}000 steps, minibatch 256, 15 SGD epochs per
batch, rollout fragments of 200 steps. The entropy coefficient follows a
schedule from 0.03 to 0.005 with the transition anchored at 83\% of the
total budget (step 2.49M of 3M) --- expressed as a budget fraction so that
short diagnostic runs and full runs traverse the same exploration profile.
The full configuration files are reproduced in \cref{app:hyperparams};
the empirical justification of the non-default values (notably $\gamma$,
the value-loss settings, and the entropy schedule) is part of the
iterative record in \cref{ch:methodology}.

\subsection{Evaluation pipeline}
\label{sec:eval-pipeline}

Results are measured at three levels, all computed from agent-level
records (\cref{sec:termination}):

\begin{enumerate}
  \item \textbf{Training stream}: the cumulative metrics over all training
        episodes. These are on-policy and, in the curriculum arm, computed
        over a level mix that changes over time --- they characterize the
        training process, not final ability.
  \item \textbf{Static recomputation}: an offline pass over the training
        logs that recomputes all aggregates from disk, deduplicates
        records, and verifies integrity (exactly six records per episode,
        no non-finite values).
  \item \textbf{Final checkpoint evaluation}: the trained checkpoint is
        reloaded and run on four fixed scenarios of 100 episodes each ---
        easy, medium and hard on Town03, plus the test level on Town05
        (\cref{sec:routes}), the latter measuring map generalization on
        roads never seen in training. Episodes are evaluated with the
        stochastic policy (sampling from the learned Gaussian) and
        per-episode reseeding of both the scenario and the action sampler.
        Evaluating the deterministic policy mean was deliberately
        rejected: for a continuous Gaussian policy it collapses behaviour
        onto a single trajectory per scenario, a failure mode analysed in
        \cref{sec:bugs}.
\end{enumerate}
